%%%%%%%%%%%%%%%%%%%%%%%%%%%%%%%%%%%%%%%%
% datoteka diploma-FRI-vzorec.tex
%
% vzorčna datoteka za pisanje diplomskega dela v formatu LaTeX
% na UL Fakulteti za računalništvo in informatiko
%
% na osnovi starejših verzij vkup spravil Franc Solina, maj 2021
% prvo verzijo je leta 2010 pripravil Gašper Fijavž
%
% za upravljanje z literaturo ta vezija uporablja BibLaTeX
%
% svetujemo uporabo Overleaf.com - na tej spletni implementaciji LaTeXa ta vzorec zagotovo pravilno deluje
%

\documentclass[a4paper,12pt,openright]{book}
%\documentclass[a4paper, 12pt, openright, draft]{book}  Nalogo preverite tudi z opcijo draft, ki pokaže, katere vrstice so predolge! Pozor, v draft opciji, se slike ne pokažejo!
 
\usepackage[utf8]{inputenc}   % omogoča uporabo slovenskih črk kodiranih v formatu UTF-8
\usepackage[slovene,english]{babel}    % naloži, med drugim, slovenske delilne vzorce
\usepackage[pdftex]{graphicx}  % omogoča vlaganje slik različnih formatov
\usepackage{fancyhdr}          % poskrbi, na primer, za glave strani
\usepackage{amssymb}           % dodatni matematični simboli
\usepackage{amsmath}           % eqref, npr.
\usepackage{hyperxmp}
\usepackage{mathtools}
\usepackage{todonotes}

\usepackage[hyphens]{url}
\usepackage{csquotes}
\usepackage[pdftex, colorlinks=true,
						citecolor=black, filecolor=black, 
						linkcolor=black, urlcolor=black,
						pdfproducer={LaTeX}, pdfcreator={LaTeX}]{hyperref}

\usepackage{color}
\usepackage{soul}

\usepackage[
backend=biber,
style=numeric,
sorting=nty,
]{biblatex}


\addbibresource{literatura.bib} %Imports bibliography file


%%%%%%%%%%%%%%%%%%%%%%%%%%%%%%%%%%%%%%%%
%	DIPLOMA INFO
%%%%%%%%%%%%%%%%%%%%%%%%%%%%%%%%%%%%%%%%
\newcommand{\ttitle}{Vzorec diplomskega dela}
\newcommand{\ttitleEn}{Diploma thesis template}
\newcommand{\tsubject}{\ttitle}
\newcommand{\tsubjectEn}{\ttitleEn}
\newcommand{\tauthor}{Matej Klančar}
\newcommand{\tkeywords}{računalnik, računalnik, računalnik}
\newcommand{\tkeywordsEn}{computer, computer, computer}
\newcommand{\proj}{\mathcal{P}}
\newcommand{\shrink}{\mathcal{D}}


%%%%%%%%%%%%%%%%%%%%%%%%%%%%%%%%%%%%%%%%
%	HYPERREF SETUP
%%%%%%%%%%%%%%%%%%%%%%%%%%%%%%%%%%%%%%%%
\hypersetup{pdftitle={\ttitle}}
\hypersetup{pdfsubject=\ttitleEn}
\hypersetup{pdfauthor={\tauthor}}
\hypersetup{pdfkeywords=\tkeywordsEn}

%%%%%%%%%%%%%%%%%%%%%%%%%%%%%%%%%%%%%%%%
% postavitev strani
%%%%%%%%%%%%%%%%%%%%%%%%%%%%%%%%%%%%%%%%  

\addtolength{\marginparwidth}{-20pt} % robovi za tisk
\addtolength{\oddsidemargin}{40pt}
\addtolength{\evensidemargin}{-40pt}

\renewcommand{\baselinestretch}{1.3} % ustrezen razmik med vrsticami
\setlength{\headheight}{15pt}        % potreben prostor na vrhu
\renewcommand{\chaptermark}[1]%
{\markboth{\MakeUppercase{\thechapter.\ #1}}{}} \renewcommand{\sectionmark}[1]%
{\markright{\MakeUppercase{\thesection.\ #1}}} \renewcommand{\headrulewidth}{0.5pt} \renewcommand{\footrulewidth}{0pt}
\fancyhf{}
\fancyhead[LE,RO]{\sl \thepage} 
%\fancyhead[LO]{\sl \rightmark} \fancyhead[RE]{\sl \leftmark}
\fancyhead[RE]{\sc \tauthor}              % dodal Solina
\fancyhead[LO]{\sc Diplomska naloga}     % dodal Solina


\newcommand{\BibLaTeX}{{\sc Bib}\LaTeX}
\newcommand{\BibTeX}{{\sc Bib}\TeX}

%%%%%%%%%%%%%%%%%%%%%%%%%%%%%%%%%%%%%%%%
% naslovi
%%%%%%%%%%%%%%%%%%%%%%%%%%%%%%%%%%%%%%%%  

\newcommand{\autfont}{\Large}
\newcommand{\titfont}{\LARGE\bf}
\newcommand{\clearemptydoublepage}{\newpage{\pagestyle{empty}\cleardoublepage}}
\setcounter{tocdepth}{1}	      % globina kazala

%%%%%%%%%%%%%%%%%%%%%%%%%%%%%%%%%%%%%%%%
% konstrukti
%%%%%%%%%%%%%%%%%%%%%%%%%%%%%%%%%%%%%%%%  
\newtheorem{izrek}{Izrek}[chapter]
\newtheorem{trditev}{Trditev}[izrek]
\newenvironment{dokaz}{\emph{Dokaz.}\ }{\hspace{\fill}{$\Box$}}


%%%%%%%%%%%%%%%%%%%%%%%%%%%%%%%%%%%%%%%%%%%%%%%%%%%%%%%%%%%%%%%%%%%%%%%%%%%%%%%
%% PDF-A
%%%%%%%%%%%%%%%%%%%%%%%%%%%%%%%%%%%%%%%%%%%%%%%%%%%%%%%%%%%%%%%%%%%%%%%%%%%%%%%

%%%%%%%%%%%%%%%%%%%%%%%%%%%%%%%%%%%%%%%% 
% define medatata
%%%%%%%%%%%%%%%%%%%%%%%%%%%%%%%%%%%%%%%% 
\def\Title{\ttitle}
\def\Author{\tauthor, matjaz.kralj@fri.uni-lj.si}
\def\Subject{\ttitleEn}
\def\Keywords{\tkeywordsEn}

%%%%%%%%%%%%%%%%%%%%%%%%%%%%%%%%%%%%%%%% 
% \convertDate converts D:20080419103507+02'00' to 2008-04-19T10:35:07+02:00
%%%%%%%%%%%%%%%%%%%%%%%%%%%%%%%%%%%%%%%% 
\def\convertDate{%
    \getYear
}

{\catcode`\D=12
 \gdef\getYear D:#1#2#3#4{\edef\xYear{#1#2#3#4}\getMonth}
}
\def\getMonth#1#2{\edef\xMonth{#1#2}\getDay}
\def\getDay#1#2{\edef\xDay{#1#2}\getHour}
\def\getHour#1#2{\edef\xHour{#1#2}\getMin}
\def\getMin#1#2{\edef\xMin{#1#2}\getSec}
\def\getSec#1#2{\edef\xSec{#1#2}\getTZh}
\def\getTZh +#1#2{\edef\xTZh{#1#2}\getTZm}
\def\getTZm '#1#2'{%
    \edef\xTZm{#1#2}%
    \edef\convDate{\xYear-\xMonth-\xDay T\xHour:\xMin:\xSec+\xTZh:\xTZm}%
}

%\expandafter\convertDate\pdfcreationdate 

%%%%%%%%%%%%%%%%%%%%%%%%%%%%%%%%%%%%%%%%
% get pdftex version string
%%%%%%%%%%%%%%%%%%%%%%%%%%%%%%%%%%%%%%%% 
\newcount\countA
\countA=\pdftexversion
\advance \countA by -100
\def\pdftexVersionStr{pdfTeX-1.\the\countA.\pdftexrevision}


%%%%%%%%%%%%%%%%%%%%%%%%%%%%%%%%%%%%%%%%
% XMP data
%%%%%%%%%%%%%%%%%%%%%%%%%%%%%%%%%%%%%%%%  
\usepackage{xmpincl}
%\includexmp{pdfa-1b}

%%%%%%%%%%%%%%%%%%%%%%%%%%%%%%%%%%%%%%%%
% pdfInfo
%%%%%%%%%%%%%%%%%%%%%%%%%%%%%%%%%%%%%%%%  
\pdfinfo{%
    /Title    (\ttitle)
    /Author   (\tauthor, damjan@cvetan.si)
    /Subject  (\ttitleEn)
    /Keywords (\tkeywordsEn)
    /ModDate  (\pdfcreationdate)
    /Trapped  /False
}

%%%%%%%%%%%%%%%%%%%%%%%%%%%%%%%%%%%%%%%%
% znaki za copyright stran
%%%%%%%%%%%%%%%%%%%%%%%%%%%%%%%%%%%%%%%%  

\newcommand{\CcImageCc}[1]{%
	\includegraphics[scale=#1]{cc_cc_30.pdf}%
}
\newcommand{\CcImageBy}[1]{%
	\includegraphics[scale=#1]{cc_by_30.pdf}%
}
\newcommand{\CcImageSa}[1]{%
	\includegraphics[scale=#1]{cc_sa_30.pdf}%
}

%%%%%%%%%%%%%%%%%%%%%%%%%%%%%%%%%%%%%%%%%%%%%%%%%%%%%%%%%%%%%%%%%%%%%%%%%%%%%%%
%%%%%%%%%%%%%%%%%%%%%%%%%%%%%%%%%%%%%%%%%%%%%%%%%%%%%%%%%%%%%%%%%%%%%%%%%%%%%%%

\begin{document}
\selectlanguage{slovene}
\frontmatter
\setcounter{page}{1} %
\renewcommand{\thepage}{}       % preprečimo težave s številkami strani v kazalu

%%%%%%%%%%%%%%%%%%%%%%%%%%%%%%%%%%%%%%%%
%naslovnica
 \thispagestyle{empty}%
   \begin{center}
    {\large\sc Univerza v Ljubljani\\%
%      Fakulteta za elektrotehniko\\% za študijski program Multimedija
%      Fakulteta za upravo\\% za študijski program Upravna informatika
      Fakulteta za računalništvo in informatiko\\%
%      Fakulteta za matematiko in fiziko\\% za študijski program Računalništvo in matematika
     }
    \vskip 10em%
    {\autfont \tauthor\par}%
    {\titfont \ttitle \par}%
    {\vskip 3em \textsc{DIPLOMSKO DELO\\[5mm]         % dodal Solina za ostale študijske programe
%    VISOKOŠOLSKI STROKOVNI ŠTUDIJSKI PROGRAM\\ PRVE STOPNJE\\ RAČUNALNIŠTVO IN INFORMATIKA}\par}%
     UNIVERZITETNI  ŠTUDIJSKI PROGRAM\\ PRVE STOPNJE\\ RAČUNALNIŠTVO IN INFORMATIKA}\par}%
%    INTERDISCIPLINARNI UNIVERZITETNI\\ ŠTUDIJSKI PROGRAM PRVE STOPNJE\\ MULTIMEDIJA}\par}%
%    INTERDISCIPLINARNI UNIVERZITETNI\\ ŠTUDIJSKI PROGRAM PRVE STOPNJE\\ UPRAVNA INFORMATIKA}\par}%
%    INTERDISCIPLINARNI UNIVERZITETNI\\ ŠTUDIJSKI PROGRAM PRVE STOPNJE\\ RAČUNALNIŠTVO IN MATEMATIKA}\par}%
    \vfill\null%
% izberite pravi habilitacijski naziv mentorja!
    {\large \textsc{Mentor}: doc.\ dr.\ 
    Aljaž Zalar\par}%
    {\vskip 2em \large Ljubljana, \the\year \par}%
\end{center}
% prazna stran
%\clearemptydoublepage      
% izjava o licencah itd. se izpiše na hrbtni strani naslovnice

%%%%%%%%%%%%%%%%%%%%%%%%%%%%%%%%%%%%%%%%
%copyright stran
%%%%%%%%%%%%%%%%%%%%%%%%%%%%%%%%%%%%%%%%
\newpage
\thispagestyle{empty}

\vspace*{5cm}
{\small \noindent
To delo je ponujeno pod licenco \textit{Creative Commons Priznanje avtorstva-Deljenje pod enakimi pogoji 2.5 Slovenija} (ali novej\v so razli\v cico).
To pomeni, da se tako besedilo, slike, grafi in druge sestavine dela kot tudi rezultati diplomskega dela lahko prosto distribuirajo,
reproducirajo, uporabljajo, priobčujejo javnosti in predelujejo, pod pogojem, da se jasno in vidno navede avtorja in naslov tega
dela in da se v primeru spremembe, preoblikovanja ali uporabe tega dela v svojem delu, lahko distribuira predelava le pod
licenco, ki je enaka tej.
Podrobnosti licence so dostopne na spletni strani \href{http://creativecommons.si}{creativecommons.si} ali na Inštitutu za
intelektualno lastnino, Streliška 1, 1000 Ljubljana.

\vspace*{1cm}
\begin{center}% 0.66 / 0.89 = 0.741573033707865
\CcImageCc{0.741573033707865}\hspace*{1ex}\CcImageBy{1}\hspace*{1ex}\CcImageSa{1}%
\end{center}
}

\vspace*{1cm}
{\small \noindent
Izvorna koda diplomskega dela, njeni rezultati in v ta namen razvita programska oprema je ponujena pod licenco GNU General Public License,
različica 3 (ali novejša). To pomeni, da se lahko prosto distribuira in/ali predeluje pod njenimi pogoji.
Podrobnosti licence so dostopne na spletni strani \url{http://www.gnu.org/licenses/}.
}

\vfill
\begin{center} 
\ \\ \vfill
{\em
Besedilo je oblikovano z urejevalnikom besedil \LaTeX.}
\end{center}

% prazna stran
\clearemptydoublepage

%%%%%%%%%%%%%%%%%%%%%%%%%%%%%%%%%%%%%%%%
% stran 3 med uvodnimi listi
\thispagestyle{empty}
\
\vfill

\bigskip
\noindent\textbf{Kandidat:} Matej Klančar\\
\noindent\textbf{Naslov:} Naslov diplomskega dela\\
% vstavite ustrezen naziv študijskega programa!
\noindent\textbf{Vrsta naloge:} Diplomska naloga na univerzitetnem programu prve stopnje Računalništvo in informatika \\
% izberite pravi habilitacijski naziv mentorja!
\noindent\textbf{Mentor:} doc.\ dr.\ Aljaž Zalar\\

\bigskip
\noindent\textbf{Opis:}\\
Besedilo teme diplomskega dela študent prepiše iz študijskega informacijskega sistema, kamor ga je vnesel mentor. 
V nekaj stavkih bo opisal, kaj pričakuje od kandidatovega diplomskega dela. 
Kaj so cilji, kakšne metode naj uporabi, morda bo zapisal tudi ključno literaturo.

\bigskip
\noindent\textbf{Title:} Naslov diplomskega dela v angleščini

\bigskip
\noindent\textbf{Description:}\\
opis diplome v angleščini

\vfill



\vspace{2cm}

% prazna stran
\clearemptydoublepage

% zahvala
\thispagestyle{empty}\mbox{}\vfill\null\it%
\noindent
Na tem mestu zapišite, komu se zahvaljujete za pomoč pri izdelavi diplomske naloge oziroma pri vašem študiju nasploh. Pazite, da ne boste koga pozabili. Utegnil vam bo zameriti. Temu se da izogniti tako, da celotno zahvalo izpustite.
\rm\normalfont

% prazna stran
\clearemptydoublepage

%%%%%%%%%%%%%%%%%%%%%%%%%%%%%%%%%%%%%%%%
% posvetilo, če sama zahvala ne zadošča :-)
\thispagestyle{empty}\mbox{}{\vskip0.20\textheight}\mbox{}\hfill\begin{minipage}{0.55\textwidth}%
Svoji dragi Alenčici.
\normalfont\end{minipage}

% prazna stran
\clearemptydoublepage


%%%%%%%%%%%%%%%%%%%%%%%%%%%%%%%%%%%%%%%%
% kazalo
\pagestyle{empty}
\def\thepage{}% preprečimo težave s številkami strani v kazalu
\tableofcontents{}


% prazna stran
\clearemptydoublepage

%%%%%%%%%%%%%%%%%%%%%%%%%%%%%%%%%%%%%%%%
% seznam kratic

\chapter*{Seznam uporabljenih kratic}

\noindent\begin{tabular}{p{0.11\textwidth}|p{.39\textwidth}|p{.39\textwidth}}    % po potrebi razširi prvo kolono tabele na račun drugih dveh!
  {\bf kratica} & {\bf angleško}                              & {\bf slovensko} \\ \hline
  {\bf CA}      & classification accuracy               & klasifikacijska točnost \\
  {\bf DBMS} & database management system & sistem za upravljanje podatkovnih baz \\
  {\bf SVM}   & support vector machine              & metoda podpornih vektorjev \\
%  \dots & \dots & \dots \\
\end{tabular}


% prazna stran
\clearemptydoublepage

%%%%%%%%%%%%%%%%%%%%%%%%%%%%%%%%%%%%%%%%
% povzetek
\phantomsection
\addcontentsline{toc}{chapter}{Povzetek}
\chapter*{Povzetek}

\noindent\textbf{Naslov:} \ttitle
\bigskip

\noindent\textbf{Avtor:} \tauthor
\bigskip

%\noindent\textbf{Povzetek:} 
\noindent V vzorcu je predstavljen postopek priprave diplomskega dela z uporabo okolja \LaTeX. Vaš povzetek mora sicer vsebovati približno 100 besed, ta tukaj je odločno prekratek.
Dober povzetek vključuje: (1) kratek opis obravnavanega problema, (2) kratek opis vašega pristopa za reševanje tega problema in (3) (najbolj uspešen) rezultat ali prispevek diplomske naloge.

\bigskip

\noindent\textbf{Ključne besede:} \tkeywords.
% prazna stran
\clearemptydoublepage

%%%%%%%%%%%%%%%%%%%%%%%%%%%%%%%%%%%%%%%%
% abstract
\phantomsection
\selectlanguage{english}
\addcontentsline{toc}{chapter}{Abstract}
\chapter*{Abstract}

\noindent\textbf{Title:} \ttitleEn
\bigskip

\noindent\textbf{Author:} \tauthor
\bigskip

%\noindent\textbf{Abstract:} 
\noindent This sample document presents an approach to typesetting your BSc thesis using \LaTeX. 
A proper abstract should contain around 100 words which makes this one way too short.
\bigskip

\noindent\textbf{Keywords:} \tkeywordsEn.
\selectlanguage{slovene}
% prazna stran
\clearemptydoublepage

%%%%%%%%%%%%%%%%%%%%%%%%%%%%%%%%%%%%%%%%
\mainmatter
\setcounter{page}{1}
\pagestyle{fancy}

\chapter{Uvod}

% Optional TOC
% \tableofcontents
% \pagebreak
\chapter{Pregled sorodnih del}


%--Paper--
\cite{BK07} Opisuje tehnike, ki jih je uporabila zmagovalna ekipa Netflixove nagrade.
Poglavje 3.1 opisuje uporabo k-NN metode ter izračun uteži glede na podatke o uporabnikih in filmih.

Poglavje 3.2 opisuje modele latentnih faktorjev. Metoda uporablja regulitacijo 
s členom $\lambda (||p_u||^2 + ||q_i||^2)$, kjer so ugotovili,
da je najboljše določiti $\lambda=0.05$. Algoritem omogoča iskanje medsebojnih vplivov
uporabnikov in filmov zaradi določenih lastnosti.

\cite{CR08} Opisuje uporabo matričnih napolnitev.
Poglavje 1.1.1 opisuje katere matrike lahko napolnimo.
Poglavja 1.1.3, 1.3 opisujejo katere algoritme je vredno obravnavati, poglavje 1.4 pa kako lahko to dosežemo z uporabo semi-definitnega programiranja.

Poglavje 2 opisuje inkoherenco, vrednost, ki nam pomaga določiti verjetnost,
da bomo matriko lahko obnovili.

Poglavje 7 definira preizkus uporabe algoritmov matričnih napolnitev,
ter predstavlja rezultate kdaj je bilo matrike možno obnoviti in kdaj ne.

\cite{VV20} Kratek video o matričnih napolnitvah, ki vsebuje preprost a počasen algoritem 
za reševanje problema.

\section{Singularni razcep}

Pri iskanju singularnega razcepa (SVD) matrike $A \in \mathbb{R}^{n_1 \times n_2}$ iščemo tri matrike, tako da velja

\[
A = U \Sigma V^T
\]

kjer velja, da je matrika $U \in \mathbb{R}^{n_1 \times n_1}$ ortogonalna in sestavljena iz levih singularnih vektorjev. Matrika $\Sigma \in \mathbb{R}^{n_1 \times n_2}$ je diagonalna, elementi na diagonali pa so po velikosti urejene singularne vrednosti matrike $A$. Matrika $V \in \mathbb{R}^{n_2 \times n_2}$ je ortogonalna, vendar sestavljena iz desnih singularnih vektorjev. 

Razcep SVD je bližnje povezan z dekompozicijo lastnih vrednosti simetričnih matrik $A^T A$, $A A^T$ ter 
$\begin{bsmallmatrix}
  0 & A^T\\
  A & 0
\end{bsmallmatrix}$.
Algoritmi nato reducirajo matriko $A$  na $A = U_1BV_1^T$, kjer sta $U_1$ ter $V_1^T$ ortogonalna, B pa ima elemente le na glavni diagonali ter na diagonali nad njo. Nato poiščemo SVD razcep matrike $B = U_2 \Sigma V_2^T$, ter razcep vstavimo v prejšnjo enačbo. Tako dobimo $A = U_1U_2 \Sigma (V_1V_2)^T$. Sedaj lahko vidimo da velja $U = U_1U_2$ ter $V = V_1V_2$.\cite{demmel97}
\chapter{Algoritmi}

\CB{Na koncu napisati kratek odstavek o tem, kaj je vsebina poglavja.}

\section{Pomembe definicije}
Nekatere definicije so uporabljene čez več algoritmov. Z namenom preglednosti, te opisujem v tem poglavju
\begin{enumerate}
  \item $\Omega$ je definirana množica znanih vrednosti
  \item \[ [\proj]_{i,j} = \begin{cases}
            a_{ij} & (i, j) \in \Omega \\
            0      & \textrm{drugače}
          \end{cases}
        \]
  \item Operator $\shrink_\tau$ kot \[
            \shrink_\tau(A) := U \shrink_\tau(\Sigma) V^T, \hspace{0.3cm} \shrink_\tau(\Sigma) = diag(max(\sigma_i - \tau, 0))
        \] \cite{CCS}
  \item Z oznako $M \in \mathbb{R}^{n_1 \times n_2}$ označujemo vhodno matriko, torej tisto, ki ima nekatere podatke neznane.
  % \item Z oznako $M \in \mathbb{R}^{n_1 \times n_2}$ označujemo bitno matriko M, ki predstavlja masko, s katero označimo katere vrednosti poznamo in katere ne. Vrednost 1 označuje, da je vrednost na tisti poziciji znana, medtem ko 0 označuje, da ni. 
\end{enumerate}

\section{Minimizacija nuklearne norme}
Ker je minimizacija ranga matrike NP-poln problem, so se razvile druge metode, ki samo kompleksnost problema zmanjšujejo. Minimizacija nuklearne norme (Nuclear Norm Minimization) se zanaša na dejstvo, da je rang matrike povezan z nuklearno normo matrike. Ta je definirana kot
\[
  ||X||_* = \sum_{i = i}^{n} \sigma_i(X)
\] \todo{se lahko sklicujem?}
Dokazano je bilo, da nuklearna norma predstavlja konveksno ovojnico ranga. \cite{NNM-PHD} Konveksna ovojnica funkcije $f : \mathcal{C} \rightarrow \mathbb{R}$ je največja konveksna funkcija $g$ tako da velja $f(x) \geq g(x)$ za vse $x \in \mathcal{C}$. \cite{Survey-NKS19}

Minimizacijo nuklearne norme je možno pretvoriti v semidefinitni problem, ki ga lahko rešujemo z različnimi pripomočki, na primer SeDuMi \cite{SeDuMi}.

Standardna oblika semidefinitnega programa je
\begin{align*}
    \min_Y \hspace{0.5cm} &\trOp{C}{Y} \\
    \text{tako da} \hspace{0.5cm} &\trOp{A_k}{Y} = b_k, \hspace{0.5cm} k = 1, \hdots , l\\
    &Y \succeq 0
\end{align*}
Kjer je $C$ dana matrika, $\{A_k\}$ in $\{b_k\}$ pa množici matrik in vektorjev. 

Problem minimizacije nuklearne norme lahko zapišemo kot 
\begin{align*}
    \min_{X, t} \hspace{0.5cm} &t \\
    \text{tako da} \hspace{0.5cm} &||X||_* \leq t,\\
    &\proj(X) = \proj(M)
\end{align*}
Trdimo, da za matriko $X \in \mathbb{R}^{n_1 \times n_2}$ in $t \in \mathbb{R}$ 
velja $||X||_* \leq t \iff \exists W_1\in \mathbb{R}^{n_1 \times n_1}, W_2 \in \mathbb{R}^{n_2 \times n_2}$ tako da \cite{NNM-PHD} \todo{se lahko sklicujem?}
\begin{align*}
    &Y = \begin{bmatrix}
        W_1 & X \\
        X^T & W_2
    \end{bmatrix}, \\
    &Y \succeq 0, \hspace{0.2cm} \tr(Y) \leq 2t
\end{align*} 

Minimizacijski problem lahko tako redefiniramo kot 
\begin{align*}
    \min_{Y,t} \hspace{0.5cm} &2t\\
    \text{tako da} \hspace{0.5cm} & Y \succeq 0,\\
    &\trOp{Y}{A_{a,b}} = b_{a,b}
\end{align*} 
kjer je $A_{a,b} \in \mathbb{R}^{(n_1 + n_2) \times (n_1 + n_2)}$ matrika v množici $A$. Velja 
$A_{a, b} \in A \iff (a, b) \in \Omega$. Matrika $A_{a,b}$ ima vse elemente ničelne, razen na mestu $(a, n_1 + b)$, kjer ima vrednost $1$.
Podobno je definiran tudi $b_{a,b} \in \mathbb{R}$, ki ima vrednost $M_{a, b}$.
Ker velja 
\[
    \trOp{A}{B} = \sum_{i}^{n_1} \sum_{j}^{n_2} a_{i,j}b_{i,j}
\] je lahko videti, da je tak pogoj smiselen. Programi za reševanje SDP pa lahko tako obliko že sprejmejo.
\cite{Survey-NKS19}


% Za SVD razcep matrike $X = U\Sigma V^T$ velja 
% \[
%     \tr \begin{bmatrix}
%         UU^T & -UV^T \\
%         -VU^T & VV^T
%     \end{bmatrix}
%     \begin{bmatrix}
%         Y & X \\
%         X^T & Z
%     \end{bmatrix} \geq 0
% \] \todo{zakaj je prva matrika pozitivno semidefinitna?}
% ker je vsota diagonalnih elementov produkta dveh pozitivnih semidefinitnih matrik vedno nenegativna. Tako vemo, da 
% \[
%     \tr(UU^TY) - \tr(UV^TX^T) - \tr(VU^TX) + \tr(VV^TZ) \geq 0
% \]

% Po \cite{CR08} lahko problem definiramo kot
% \begin{align*}
%   \min    & \hspace{0.5cm} \tr(Y)                                    \\
%   \text{tako da} & \hspace{0.5cm} (Y, A_k) = b_k, k = 1, \cdots , |\Omega| \\
%                     & \hspace{0.5cm} Y \succcurlyeq 0
% \end{align*}
% kjer
% \begin{align*}
%   &Y = \begin{bmatrix}
%     W_1 & X   \\
%     X^T & W_2
%   \end{bmatrix}
% \end{align*} tak problem pa lahko že rešujemo s semidefinitnimi programi.
\section{Algoritem praga singularnih vrednosti}\label{2807-1441}
\textbf{Algoritem praga singularnih vrednosti (SVT)} \cite{CCS} uporabi idejo, da imajo matrike z majhnim rangom nekaj velikih singularnih vrednosti, ostale pa 0 ali pa vsaj blizu 0. Ključna parametra v SVT-ju sta \textit{izbira premika} in \textit{izbira praga},  
%Za svoje delovanje uvede dva nova pomembna koncepta, prvi je premik, drugi pa prag, potreben za uporabo operatorja $\shrink_\tau$ \eqref{1007-1959}. 
algoritem pa temelji na iteraciji
\begin{align}
\label{2407-1910}
        X^{(k)} &= \shrink_\tau(Y^{(k-1)}), \\
        Y^{(k)} &= Y^{(k-1)} + \delta_k \proj(M - X^{(k)}), 
\end{align}
kjer so $\tau > 0$ izbran prag, $\delta_k$ izbran premik, $X^{(0)} = 0 \in \mathbb{R}^{n_1 \times n_2}$ in
$Y^{(0)} = 0 \in \mathbb{R}^{n_1 \times n_2}$. \cite{CCS}

V nadaljevanju bomo opisali glavno idejo zgornje iteracije. V grobem pa temelji na uporabi metode za iskanje vezanih ekstremov, kjer elementi matrik $Y^{(k)}$ predstavljajo Lagrangove množitelje. 

Uvedimo funkcijo 
\begin{align}
    \label{1007-2007}
    f_\tau(X) = \tau\nnorm{X} + \frac{1}{2}\fnorm{X}^2
\end{align}
in optimizacijski problem
\begin{align}
\label{2706-0957}
\begin{split}
    \min_{X\in \mathbb R^{n_1\times n_2}} & \hspace{0.5cm} f_\tau(X), \\
    \textrm{pri pogojih} & \hspace{0.5cm} \proj(X) = \proj(M).
\end{split}
\end{align}
Opazimo lahko, da za velike vrednosti $\tau$ velja $f_\tau(X) \approx \tau\nnorm{X}$, kar pomeni, da bo s primerno izbranim $\tau$, optimizacijski problem minimiziral nuklearno normo.

Denimo, da želimo poiskati minimum funkcije $f(x)$ pri pogojih $g_1(x)=g_2(x)=\ldots=g_{k}(x)=0$.
V teoriji vezanih ekstremov se za tovrstne probleme uvede \textbf{Lagrangeovo funkcijo}
\[\mathcal{L}(x, \lambda_1,\lambda_2,\ldots,\lambda_k) = f(x) + \lambda_1 g_1(x)+\lambda_2g_2(x)+\ldots+\lambda_kg_k(x),\]
nato pa išče ekstreme med njenimi stacionarnimi točkami.
Problemu \eqref{2706-0957} lahko priredimo Lagrangeovo funkcijo
\[
    \mathcal{L}(X, Y) = f_\tau(X) + \left< Y, \proj(M - X) \right>,
\] 
nato pa iščemo njene stacionarne točke.
Zaradi velikega števila parametrov pa ta pristop navadno ni izvedljiv, zato se v SVT-ju za iskanje ekstremov $\mathcal{L}(X, Y)$ uporabi t.i.\ \textit{Uzawa algoritem} \cite{CCS}. Ta ekstreme išče prek iterativnega postopka:
\begin{align}
        X^{(k)} &=  \argmin_{X} \mathcal{L}(X^{(k)}, Y^{(k-1)}) \label{1007-2018},\\
        Y^{(k)} &= Y^{(k-1)} + \delta_k \proj(M - X^{(k)}). \label{1007-2019}
\end{align}

Izkaže se, da je rešitev \eqref{1007-2018} enaka $\shrink_\tau(Y^{(k-1)})$. To spodaj dokažemo, še prej pa izpeljimo pomožen rezultat.

\begin{trditev}
Velja:
\begin{align}
    \argmin_X  \mathcal{L}(X, Y)
    = \argmin_X \Big(\tau\nnorm{X} + \frac{1}{2}\fnorm{X - Y}^2\Big)
\end{align}
\end{trditev}

\begin{proof}
Trditev sledi iz krajšega računa:
\begin{align*}
    &\argmin_X  \Big(\tau\nnorm{X} + \frac{1}{2}\fnorm{X - Y}^2\Big) \\
    &=\argmin_X \Big(\tau\nnorm{X} + \frac{1}{2}\trOp{X - Y}{X - Y}\Big)\\
    &=\argmin_X \Big(\tau\nnorm{X} + \frac{1}{2}(\fnorm{X}^2 - 2\trOp{X}{Y} + \fnorm{Y}^2)\Big) \\ 
    &=\argmin_X \Big(\tau\nnorm{X} + \frac{1}{2}\fnorm{X}^2 - \trOp{X} {\proj(Y)}\Big)\\
    &=\argmin_X \Big(\tau\nnorm{X} + \frac{1}{2}\fnorm{X}^2 + \tr(- \proj(X)Y^T) + \tr(\proj(M)Y^T)\Big)\\
    &=\argmin_X \Big(\tau\nnorm{X} + \frac{1}{2}\fnorm{X}^2 + \tr(\proj(M - X)Y^T)\Big)\\
    &=\argmin_X \Big(\tau\nnorm{X} + \frac{1}{2}\fnorm{X}^2 + \trOp{Y}{\proj(M - X)}\Big)\\ 
    &= \argmin_X \mathcal{L}(X, Y)
\end{align*}
kjer smo v prvi enakosti uporabili definicijo Frobeniusove norme, v drugi bilinearnost skalarnega produkta, 
v tretji smo ignorirali konstanto $\fnorm{Y}^2$, saj ne vpliva na rezultat, 
in upoštevali $\proj(Y^{(k)}) = Y^{(k)}$ za vse $k \in \mathbb{N}$. Zadnje dejstvo sledi iz definicije \eqref{1007-2019} in $Y^{(0)} = 0$. V četrti enakosti smo upoštevali $\trOp{X}{\proj(Y)} = \trOp{\proj(X)}{Y}$, kar je lahko videti, ko se spomnimo, da za skalarni produkt $\trOp{A}{B}$ matrik $A, B \in \mathbb{R}^{n_1 \times n_2}$ velja 
\[
    \trOp{A}{B} = \sum_{i = 1}^{n_1}\sum_{j = 1}^{n_2} a_{ij}b_{ij}.
\]
Prišteli smo tudi konstanto $\tr(\proj(M)Y^T)$, ki ne vpliva na rezultat.
V peti enakosti smo upoštevali linearnost sledi in aditivnost operatorja $P_\Omega$,
v šesti definicijo skalarnega produkta in v zadnji definicijo Lagrangeove funkcije.
\end{proof}
\begin{theorem} \label{1907-2240}
Za matriki $X \in \mathbb{R}^{n_1 \times n_2}$ in 
$Y \in \mathbb{R}^{n_1 \times n_2}$ velja:
\begin{align}
    \label{2906-1056}
    \shrink_\tau(Y) = \argmin_{X} \Big(\frac{1}{2} \fnorm{X-Y}^2 + \tau\nnorm{X} \Big) 
\end{align}
\end{theorem}

\begin{proof} 
Najprej se spomnimo definicije konveksne funkcije. Funkcija $f$ je konveksna, če za katerikoli dve točki $x_1, x_2$ v domeni funkcije $f$ velja, da je premica čez ti dve točki na odseku med tema dvema točkama nad grafom funkcije $f$.
Funkcija $h(X) := \frac{1}{2} \fnorm{X-Y}^2 + \tau\nnorm{X} $ je konveksna funkcija, saj potreben pogoj trikotniške neenakosti matričnih norm zagotavlja, da je matrična norma konveksna funkcija. Vsota konveksnih funkcij pa je prav tako konveksna funkcija. 
%\todo{ali je potrebno citirati} 
Zaradi konveksnosti funkcije $f$ je 
subgradient v vsaki točki iz domene dobro definiran.
Tega smo definirali v razdelku \ref{2607-1502}.
Matrika $Z\in \mathbb R^{n_1\times n_2}$ je subgradient funkcije $f$ v točki 
$X_0\in \mathbb R^{n_1\times n_2}$,
če velja 
\[\forall X\in \mathbb R^{n_1\times n_2}: 
f(X) \geq  f(X_0) + \trOp{Z}{X - X_0}.\]
Iz definicije subgradienta sledi, da bo imela funkcija $f$ minimum v točki $X'$ natanko tedaj,
ko bo ničelna matrika $\mathbf{0}$ eden izmed subgradientov funkcije $f$ v točki $X_0$.
\iffalse 
V minimumu funkcije $f$ bo eden od subgradientov iščemo minimum $X'$ funkcije $f$,  iščemo tako točko $X'$, da bo
\CG{eden izmed subgradientov po spremenljivki $X$} v točki $X'$ enak 0.  
\fi 

Izkaže se \cite{CCS}, da je množica subgradientov nuklearne norme v točki $X$ enaka
\begin{equation}  \label{2807-0932}  
    \partial\nnorm{X} = \{UV^* + W: W \in \mathbb{R}^{n_1 \times n_2}, U^*W = \mathbf{0}, WV = \mathbf{0}, \norm{W}_2 \leq 1 \},
\end{equation}
kjer $U \Sigma V^T$ predstavlja SVD razcep matrike $X$. S krajšim računom lahko preverimo, da za vsak par $(i,j)$ velja $\frac{\partial}{\partial X_{ij}}\fnorm{X-Y}^2 = 2(X_{ij} - Y_{ij})$. 
Če te parcialne odvode zložimo v matriko, dobimo $2(X-Y)$. Od tod sledi, da je $X-Y$ eden od subgradientov funkcije $h_1(X):=\frac{1}{2}\fnorm{X-Y}^2$ v točki $X$ \cite[3.1.3]{boyd2004convex}.
Iz teh dveh premislekov sledi, da bo $X'$ minimum $h$ natanko tedaj, ko velja 
\begin{equation} \label{2807-0940}
    \mathbf{0} \in X' - Y + \tau \partial\nnorm{X'}.
\end{equation}
Trditev izreka bo sledila, če pokažemo, da velja $X' = \shrink_\tau(Y)$. Najprej razčlenimo SVD razcep matrike $Y$ kot 
\[
    Y = U_0\Sigma_0V_0^T + U_1\Sigma_1V_1^T,
\]
kjer se $U_0, \Sigma_0$ in $V_0$ nanaša na singularne vrednosti in pripadajoče singularne vektorje večje od $\tau$, $U_1, \Sigma_1$ in $V_1$ pa tiste manjše od $\tau$. Pokazati želimo, da velja 
\[
    X' = U_0(\Sigma_0 - \tau I)V_0^T,
\] 
kar pa je po definiciji natanko enako operatorju $\shrink_\tau(Y)$. S pretvorbo \eqref{2807-0940} pridemo do zapisa
\begin{equation} \label{2807-0934}
    Y- X' \in \tau \partial \nnorm{X'}.
\end{equation}
Sedaj iščemo tako matriko $W$, da bodo zanjo držali pogoji, podani v \eqref{2807-0932}, prav tako pa bo po \eqref{2807-0934} držalo $Y - X' = \tau (U_0V_0^T + W)$. Spodaj bomo preverili, da je primerna izbira za $W$ enaka $W= \tau^{-1} U_1 \Sigma_1 V_1^T$. 
Veljavnost željene enakosti
preverimo s kratkima računoma:
\begin{align*}
    Y-X' &= U_0\Sigma_0V_0^T + U_1\Sigma_1V_1^T - U_0(\Sigma_0 - \tau I)V_0^T \\ 
    &= U_0(\Sigma_0 - \Sigma_0 + \tau I)V_0^T + U_1\Sigma_1 V_1^T  \\
    &= \tau U_0 V_0^T + U_1\Sigma_1 V_1^T
\end{align*}
in 
\begin{align*}
    \tau(U_0 V_0^T + W) &= \tau(U_0V_0^T + \tau^{-1} U_1 \Sigma_1 V_1^T)\\ 
    &= \tau U_0 V_0^T + U_1 \Sigma_1 V_1^T. 
\end{align*}
Sedaj je zgolj potrebno pokazati, da veljajo potrebne lastnosti matrike $W$.
Po sami definiciji SVD vemo, da so stolpci matrik $U$ in $V$ ortogonalni. Torej velja $U_0^TW = 0$ in $WV_0 = 0$. Ker pa ima matrika $\Sigma_1$ vse elemente manjše od $\tau$, velja tudi $\norm{W}_2 \leq 1$. $\norm{A}_2$ je namreč definirana kot največja singularna vrednost matrike $A$. S tem smo pokazali, da je $Y - X' \in \tau \partial \nnorm{X'}$.
\end{proof}
Z uporabo \eqref{1007-2018}, \eqref{1007-2019}
in izreka \ref{1907-2240}
res pridemo do iteracije
\eqref{2407-1910}, ki jo uporablja algoritem SVT.
\iffalse
Tako res pridemo 
Po trditvi lahko sedaj zapišemo algoritem \eqref{1007-2018} - \eqref{1007-2019} kot \cite{CCS}
\[
    \begin{cases}
        X^k = \shrink_\tau(Y^{k-1}) \\
        Y^k = Y^{k-1} + \delta_k \proj(M - X^k) 
    \end{cases}
\]
\fi

\subsection{Nastavljanje parametrov $\tau$ in $\delta$} \label{1907-1648}
Opazimo lahko, da algoritem SVT potrebuje dva parametra, $\tau$ in $\delta$, ki ju moramo izbrati 
že pred vstopom v algoritem.

Po priporočilih \cite{CCS}
sta primerni izbiri za $\delta$ in $\tau$
enaki 
\[
    \delta = 1.2\, \dfrac{n_1 n_2}{m}\qquad\text{in}\qquad
    \tau = 5n,
\]
pri čemer je $\tau$ naveden za kvadratne $n\times n$ matrike.
V poglavju z rezultati smo prvotno uporabljali ti dve konstanti, pri 
čemer smo zaradi
pravokotnosti $n_1\times n_2$ matrik uporabljali
$\tau = 5\frac{n_1+n_2}{2}$.
\iffalse
Medtem, ko so koraki v samem algoritmu definirani kot množica korakov, 
smo v okviru rezultatov diplomske naloge, prvotno za premik uporabljali konstanto, ter korak nastavili na 
 po priporočilih \cite{CCS}. 
 \fi
\iffalse
Prav tako članek \cite{CCS} navaja, da je za matrike velikosti $\mathbb{R}^{n \times n}$ smiselno nastaviti $\tau = 5n$, vendar sem v moji implementaciji zaradi posploševanja na nekvadratne matrike, za matrike velikosti $\mathbb{R}^{n_1 \times n_2}$ parameter nastavil na
\[
    \tau = 5\, \frac{n_1+n_2}{2}
\]
\fi
V naših testiranjih se je izkazalo, da sta taka parametra dobra za večje matrike. V razdelku \ref{1307-2251} pa bomo videli, da moramo pri manjših matrikah pogosto zmanjšati premik in povečati prag.
\section{Algoritem minimizacije prirezane nuklearne norme}
Že samo ime nam pove, da bo algoritem minimizacije prirezane nuklearne norme, oziroma TNNM podoben algoritmu NNM. Ideja tega algoritma pa je, da uporabimo dodatno informacijo $r \in \mathbb{N}$ o rangu originalne, nezašumljene matrike.

Sam algoritem uvede tako imenovano \textbf{$r$-prirezano nuklearno normo}, ki je za matriko $X \in \mathbb{R}^{n_1 \times n_2}$ definirana kot vsota $\min(n_1,n_2) - r$ najmanjših singularnih vrednosti
\[
    \norm{X}_r = \sum^{\min(n_1, n_2)}_{i = r + 1} \sigma_i(X)
\]
TNNM rešuje problem \cite{TNNM-HZYLH12}
\begin{align*}
    \min_X               & \hspace{0.5cm} \norm{X}_r \numberthis \label{1107-1305} \\
    \textrm{pri pogojih} & \hspace{0.5cm} \proj(X) = \proj(M)
\end{align*}
Cilj algoritma je torej čim bolj zmanjšati najmanjše singularne vrednosti, medtem ko velikih ne omejujemo. S tem problem minimizacije omilimo.

Problem \eqref{1107-1305} je ekvivalenten
\begin{align*}
    \min_X               & \hspace{0.5cm} \nnorm{X} - \sum_{i=1}^{r} \sigma_i \\
    \textrm{pri pogojih} & \hspace{0.5cm} \proj(X) = \proj(M)
\end{align*}
Za nadaljnje korake bomo potrebovali naslednji izrek
\begin{theorem}
    Za $X \in \mathbb{R}^{n_1 \times n_2}$, $A \in \mathbb{R}^{r \times n_1}$, $B \in \mathbb{R}^{r \times n_2}$ in $r \in \mathbb{N}$, pri pogojih $r \leq min(n_1, n_2)$, $AA^T = I_{r}, BB^T = I_{r}$ velja:
    \begin{align*}
        \tr(AXB^T) \leq \sum_{i=1}^{r} \sigma_i(X)
    \end{align*}
\end{theorem}

\begin{proof}
    Za dokaz uporabljamo von Neumannovo neenakost sledi \cite{TNNM-HZYLH12}, s katero lahko zapišemo
    \[
        \tr(AXB^T) = \tr(XB^TA) \leq \sum^{\min(n_1, n_2)}_{i=1} \sigma_i(X) \sigma_i(B^TA)
    \]
    Enakost $\tr(AXB^T) = \tr(XB^TA)$ sledi iz dejstva, da je sled produkta matrik invariantna pod cikličnimi permutacijami.

    Po definiciji lahko singularne vrednosti matrike $Y$ najdemo tako, da najdemo korene nenegativnih lastnih vrednosti matrike $Y^TY$. Tako lahko povemo, da so singularne vrednosti matrike $B^TA$ enake lastnim vrednostim matrike $A^TBB^TA$. Izraz razpišemo v
    \[
        A^TBB^TA = A^TI_rA = A^TA
    \]
    %https://shorturl.at/pqrRX
    Ker pa velja, da imata matriki $XY$ in $YX$ enake neničelne lastne vrednosti, ter vemo da $AA^T = I_r$, lahko povemo, da ima produkt $B^TA$ $r$ singularnih vrednosti enakih 1, saj ima $I_n$ $n$ lastnih vrednosti enakih 1.
    Tako lahko sedaj razpišemo izraz
    \[
        \sum^{\min(n_1, n_2)}_{i=1} \sigma_i(X) \sigma_i(B^TA) = \sum^{r}_{i=1} \sigma_i(X)
    \]
    Ugotovili smo, da velja
    \begin{equation}
        \label{2806-1320}
        \tr(AXB^T) \leq \sum^{r}_{i=1} \sigma_i(X)
    \end{equation}
\end{proof}

\begin{theorem}
    Za SVD razcep matrike $X = U \Sigma V^T$, ter matriki
    $A = (u_1, \hdots , u_r)^T$ in $B = (v_1, \hdots , v_r)^T$,
    kjer je $u_i$ $i$-ti stolpec matrike $U$ ter $v_i$ $i$-ti stolpec
    matrike $V$, velja \[\tr(AXB^T) = \sum^{r}_{i=1} \sigma_i(X)\]
\end{theorem}

\begin{proof}
    \begin{align*}
        \tr(AXB^T) & = \tr((u_1, \hdots , u_r)^T X (u_1, \hdots , u_r)^T)                                                                                                                                                                   \\
                   & = \tr((u_1, \hdots , u_r)^T U \Sigma V^T (u_1, \hdots , u_r)^T)                                                                                                                                                        \\
                   & = \begin{bmatrix} I_r & 0 \\ 0 & 0 \end{bmatrix} \Sigma \begin{bmatrix} I_r & 0 \\ 0 & 0 \end{bmatrix}                                                                                                                 \\
                   & = \tr(\begin{bmatrix} \sigma_1 \\[-8pt] & \ddots & & \\[-8pt] & & \sigma_r \\[-8pt] & &  & 0 \\[-8pt]  & & & & \ddots \\[-8pt] & & & & & 0 \end{bmatrix}) = \sum_{i = 1}^{r} \sigma_i(X)  \numberthis\label{2806-1321}
    \end{align*}
\end{proof}
Z združitvijo dokazov \eqref{2806-1320} in \eqref{2806-1321} zapišemo
\[
    \max_{AA^T = I, BB^T = I} \tr(AXB^T) = \sum^{r}_{i = 1} \sigma_i(X)
\]
Torej je optimizacijski problem ekvivalenten
\begin{align*}
    \min_X \hspace{0.5cm}             & \nnorm{X} - \max_{AA^T = I, BB^T = I} \tr(AXB^T) \\
    \text{pri pogojih} \hspace{0.5cm} & \proj(X) = \proj(M)
\end{align*}

Glede na vse ugotovitve, nato nastavimo iterativni algoritem, tako da, izračunamo $X^0 = \proj(M)$. V $i$-ti iteraciji izračunamo $A^i$ in $B^i$, tako da izračunamo SVD razcep $X^i = U \Sigma V^T$, ter $A$ nastavimo kot prvih $r$ stolpcev matrike $U$, $B$ pa kot prvih $r$ stolpcev matrike $V$. $X^{i+1}$ lahko sedaj izračunamo kot \cite{TNNM-HZYLH12}
\begin{align*}
    \min_X \hspace{0.5cm}         & \nnorm{X} - \tr(A^iX(B^i)^T) \\
    \text{tako da} \hspace{0.5cm} & \proj(X) = \proj(M)
\end{align*}
Če minimizacijo še malo obrnemo, pa lahko problem rešujemo z uporabo algoritma ADMM. \todo{iz kje pride izracun $Y^{i+1}$}
\begin{align*}
    \min_X \hspace{0.5cm}         & \nnorm{X} - \tr(A^iW(B^i)^T) \numberthis\label{1307-1527} \\
    \text{tako da} \hspace{0.5cm} & W=X, \hspace{0.2cm} \proj(W) = \proj(M)
\end{align*}
Ta problem ponovno zapišemo s pomočjo Lagrangeove funkcije, le da algoritem ADMM definira še regularizacijski parameter $\beta$. \cite{TNNM-HZYLH12}
\[
    \mathcal{L}(X, Y, W, \beta) = \nnorm{X} - \tr(A_l W B_l^T) + \frac{\beta}{2} \fnorm{X - W}^2 + \tr(Y^T(X-W))
\]
Opazimo lahko, da funkcija definira zgolj pogoj $X = W$. Videli bomo, da algoritem pogoj $\proj(X) = \proj(M)$ definira posredno, s popravljanjem znanih vrednosti (korak \eqref{2906-1248} spodaj).%\todo{ali lahko citiram naprej}

Matriko $X^{k+1}$ ponovno definiramo kot
\[
    X^{k+1} = \arg \min_X \mathcal{L}(X, Y^k, W^k, \beta)
\]
Z ignoriranjem konstantnih členov, pa lahko tako zapišemo
\begin{align*}
    X^{k+1} & = \arg \min_X \hspace{0.2cm} \nnorm{X} + \frac{\beta}{2}\trOp{X-W^k}{X-W^k} + \frac{\beta}{2}\trOp{\frac{2}{\beta}Y}{X}       \\
            & = \arg \min_X \hspace{0.2cm} \nnorm{X} + \frac{\beta}{2}\tr(X^TX - X^TW^k - W^{k^T}X - W^{k^T}W^k + \frac{2}{\beta}Y^{k^T}X )
\end{align*}
Z namenom faktorizacije dodamo konstantne člene.
\begin{align*}
     & \arg \min_X \hspace{0.2cm} \nnorm{X} + \frac{\beta}{2}\tr(X^T(X - W^k + \frac{1}{\beta}Y^k) - W^{k^T} (X - W^k + \frac{1}{\beta}Y^k) \\
     & +\frac{1}{\beta}Y^{k^T}(X - W^k + \frac{1}{\beta}Y^k))                                                               \\
     & = \arg \min_X \hspace{0.2cm} \nnorm{X} + \frac{\beta}{2} \fnorm{X-W^k + \frac{1}{\beta}Y^k}^2                                       \\
     & = \arg \min_X \hspace{0.2cm} \frac{1}{\beta}\nnorm{X} +  \frac{1}{2}\fnorm{X-(W^k - \frac{1}{\beta}Y^k)}^2
\end{align*}
\CB{Po enakkosti v izreku ... ali po izreku ...} Po izreku \ref{1907-2240} pa tako lahko zapišemo
\[
    X^{k+1} = \shrink_\frac{1}{\beta}(W_k - \frac{1}{\beta} Y_k)
\]
Matriko $W$ podobno izračunamo kot
\begin{align*}
    W^{k+1} & = \arg \min_{W} \mathcal{L}(X^{k+1}, Y^k, W, \beta)                                       \\
            & = \arg \min_W \frac{\beta}{2} \fnorm{W - (X^{k+1} + \frac{1}{\beta}(A_l^T B_l + Y_k)) }^2
\end{align*}
Lahko je videti, da velja
\[
    W^{k+1} = X^{k+1} + \frac{1}{\beta}(A_l^T B_l + Y_k)
\]
Zaradi pogoja $\proj(W) = \proj(M)$ pa tiste elemente, ki poznamo, popravimo. \cite{TNNM-HZYLH12}\todo{malo spremenjeno iz clanka}
\begin{equation}
    \label{2906-1248}
    W_{k+1} = W_{k+1} + \proj(M - W_{k+1})
\end{equation}
\CG{Korak preprosto spremeni mesta, kjer vrednosti poznamo na znane vrednosti, ostale pa pusti, kakršne so.}

\section{Izmenjajoč gradientni spust}
Ker je računanje SVD razcepa zahtevna operacija, saj ima časovno zahtevnost
$O(n^3)$, je bilo predlaganih nekaj algoritmov, ki za svoje delovanje ne potrebujejo SVD-ja. Algoritem Izmenjajočega gradientnega spusta, oziroma v nadaljevnu ASD sloni na računanju gradienta in premikanja po njem. Ideja algoritma je najti dve matriki $X \in \mathbb{R}^{n_1 \times r}$ ter $Y \in \mathbb{R}^{r \times n_2}$, tako da velja $\proj(M) = \proj(XY)$. Vidimo lahko, da ponovno potrebujemo informacijo o rangu matrike, ki jo rekonstruiramo. Ker imata tako $X$ in $Y$ kvečjemu rang $r$, vemo, da tudi njun produkt $XY$ ne bo imel ranga večjega od $r$. 

Cilj algoritma je minimizirati
\[
    \min_{X, Y} \hspace{0.5cm} \frac{1}{2}\, \fnorm{\proj(M) - \proj(XY)}^2
\] 
Algoritem minimizacijo razdeli na dve, nato pa izmenično rešuje eno in nato drugo kot \cite{AST-TK15}
\begin{align*}
    X_{i+1} &= \arg \min_{X} \hspace{0.3cm} \fnorm{\proj(M) - \proj(XY_i)}^2 \\
    Y_{i+1} &= \arg \min_{Y} \hspace{0.3cm} \fnorm{\proj(M) - \proj(X_{i+1}Y)}^2
\end{align*}


Za uporabo algoritma ASD potrebujemo odvoda funkcije $f(X,Y) = \frac{1}{2}\, \fnorm{\proj(M) - \proj(XY)}^2$, ki ga izračunamo za vsak element posebej.
\begin{align*}
    \frac{\partial}{\partial x_{a,b}} f &= \frac{\partial}{\partial x_{a,b}} \frac{1}{2}\, \fnorm{\proj(M) - \proj(XY)}^2  \\
    &= \frac{\partial}{\partial x_{a,b}} \frac{1}{2}\, \sum_{i}^{n_1}\sum_{j}^{n_2}(\delta_{i,j}m_{i,j} - \delta_{i,j}\sum_{k}^{r}(x_{i,k}y_{k,j}))^2  \\
    &= \, \sum_{j}^{n_2}(\delta_{a,j}m_{a,j} - \delta_{i,j}\sum_{k}^{r}(x_{a,k}y_{k,j}))(-y_{b,j}) \implies \\
    &\implies \frac{\partial}{\partial X}f = -(\proj(M) - \proj(XY))Y^T
    % &= \sum_{i,j}(\delta_{i,j}m_{i,j} - \delta_{i,j}\sum_{k}(x_{i,k}y_{k,j}))
    % (\sum_{k}y_{k,j}) = \\
    % &= (P(M) - P(XY))Y^T
\end{align*}
kjer 
\[
    \delta_{i,j} = \begin{cases}
        1, (i, j) \in \Omega \\
        0, (i, j) \notin \Omega
    \end{cases}
\]
Na podoben način bi lahko pokazali tudi
\[
    \frac{\partial}{\partial Y} = -X^T (\proj(M) - \proj(XY))
\]
Poiščimo še najboljši korak gradientnega spusta.
Cilj je pokazati, da je premik po gradientu s korakom $t_x$ najboljši.
Najprej definirajmo sam premik, kot
\[
    X^{k+1} = X^k - t_X \nabla f_Y(X)
\]
Ker želimo, da bodo znane vrednosti produkta $X^{k+1}Y^{k}$ kar se da podobne znanim vrednostim matrike $M$, nastavimo $t_x$ kot 
\begin{align*}
    t_x &= \arg \min_t \hspace{0.3cm} g(t)
\end{align*} kjer
\begin{align*}
    g(t) &= \frac{1}{2} \fnorm{\proj(M) - \proj((X - t\nabla f_Y(X))Y)}^2
\end{align*}
\CG{
\begin{align*}
    g(t) &= \frac{1}{2} \fnorm{\proj(M - XY)}^2 - t\tr(\proj(M - XY)\proj(\nabla f_Y(X)Y)^T) + \frac{t^2}{2}\fnorm{\proj(\nabla f_Y(X)Y)}^2
\end{align*}}
Ker je \cite{AST-TK15} \todo{ali lahko tako?}
\[
  g'(t) = -\fnorm{\nabla f_Y(X)}^2 + t\fnorm{\proj(f_Y(X)Y)}^2 
\]
vidimo, da funkcija $g(t)$ doseže minimum pri 
\[
  t_x = \frac{\fnorm{\nabla f_Y(X)}^2}{\fnorm{\proj(\nabla f_Y(X)Y)}^2}  
\]
Podobno velja za korak v smeri gradientnega spusta matrike $Y$, kjer
\[
  t_y = \frac{\fnorm{\nabla f_X(Y)}^2}{\fnorm{\proj(X \nabla f_X(Y))}^2}  
\]
\section{LMaFit}
LMaFit je algoritem, ki podobno kot ASD, rešuje problem matričnih napolnitev z uporabo dveh manjših matrik $X^{n_1 \times r}$ in $Y^{n_1 \times r}$. Podobno kot v prejšnjih algoritmih rešujemo problem 
\begin{align*}
    \min_{X, Y, Z}& \hspace{0.5cm} \frac{1}{2}||XY - Z||_F^2\\
    \text{tako da}& \hspace{0.5cm} \proj(M) = \proj(Z)
\end{align*}
le da sam problem rešujemo s pomočjo Moore-Penrose inverza (Označen z $\dagger$).

Lahko je videti, da bo funkcija $f(X,Y,Z) = frac{1}{2}||XY - Z||_F^2$ imela najmanjšo vrednost, ko bo $XY = Z$. Zato uvedemo iterativni algoritem, ki izmenično posodablja $X, Y$ in $Z$ tako, da fiksira dve izmed matrik. Minimizacijo matrik $X$ in $Y$ torej dosežemo na naslednji način.
\begin{align*}
    X^{i+1}Y^{i} &= Z^{i} \hspace{0.5cm} \\
    X^{i+1} &= Z^{i}(Y^i)^\dagger
\end{align*}
Kjer se zanašamo na dejstvo, da nam množenje z Moore-Penrose inverzom z desne 
da najboljši možen rezultat.

Podobno posodobimo matriko $Y$.
\begin{align*}
    \hspace{0.5cm} X^{i+1}Y^{i+1} &= Z^{i}   \\
    Y^{i+1} &= (X^{i+1})^\dagger Z^{i}
\end{align*}

Sedaj le še posodobimo matriko Z, kot
\[
    Z^{i+1} = X^{i+1}Y^{i+1} + \proj(M - X^{i+1}Y^{i+1})
\]
Torej podobno kot prej poskušamo doseči $XY = Z$, vendar zaradi omejitve znanih vrednosti matrike M, na tistem mestu vrednosti popravimo. \cite{LMaFit-WY12}





%--/Paper--
\printbibliography

\end{document}